
% !TEX program = pdflatex
\documentclass[11pt,a4paper]{article}

\usepackage[utf8]{inputenc}
\usepackage[T1]{fontenc}
\usepackage[spanish,es-nodecimaldot]{babel}
\usepackage{amsmath,amssymb,amsthm,mathtools,mathrsfs}
\usepackage{geometry}
\usepackage{enumitem}
\usepackage[expansion=false]{microtype}
\usepackage{hyperref}
\usepackage{cleveref}
\usepackage{xcolor}
\geometry{margin=2.7cm}
\setlength{\parindent}{0pt}
\setlength{\parskip}{0.6em}

\hypersetup{
	colorlinks=true,
	linkcolor=blue!60!black,
	urlcolor=blue!60!black,
	citecolor=blue!60!black
}

%------------------------------------------------
% Entornos
%------------------------------------------------
\theoremstyle{definition}
\newtheorem{definicion}{Definición}[section]
\newtheorem{ejemplo}{Ejemplo}[section]

\theoremstyle{plain}
\newtheorem{proposicion}{Proposición}[section]
\newtheorem{teorema}{Teorema}[section]
\newtheorem{corolario}{Corolario}[section]
\newtheorem{lema}{Lema}[section]

\theoremstyle{remark}
\newtheorem{observacion}{Observación}[section]

%------------------------------------------------
% Comandos
%------------------------------------------------
\newcommand{\C}{\mathbb C}
\newcommand{\R}{\mathbb R}
\newcommand{\Z}{\mathbb Z}
\newcommand{\N}{\mathbb N}
\newcommand{\I}{[0,1]}
\newcommand{\Hol}{\mathcal H}
\newcommand{\OO}{\mathcal O}
\newcommand{\id}{\operatorname{id}}
\newcommand{\restr}[2]{\left.#1\right|_{#2}}

\title{\textbf{Notas de EDO}}
\author{}
\date{}

\begin{document}
	\noindent \textbf{Exercise 41.} Prove that for all open sets $\mathcal{D}_\gamma \subseteq \mathbb{R}$ and $\mathcal{D}_f \subseteq \mathbb{C}$, and all functions $\gamma : \mathcal{D}_\gamma \to \mathbb{C}$ with a real derivative at $t \in \mathcal{D}_\gamma$, and $f : \mathcal{D}_f \to \mathbb{C}$ with a complex derivative at $z := \gamma(t) \in \mathcal{D}_f$, the composite function $f \circ \gamma$ has a real derivative at $t$ given by
	\[
	(f \circ \gamma)'(t) = \{f'[\gamma(t)]\} \cdot \gamma'(t).
	\]
	
	\subsection*{Resolución}
	
	\begin{proof}
		Queremos demostrar que la derivada real de la función compuesta $f \circ \gamma$ en el punto $t \in \mathcal{D}_\gamma$ satisface la regla de la cadena. Por definición, la derivada real de una función $\sigma : \mathcal{D}_\gamma \to \mathbb{C}$ en un punto $t$ viene dada por el límite del cociente incremental:
		\[
		\sigma'(t) = \lim_{h \to 0} \frac{\sigma(t+h) - \sigma(t)}{h}
		\]
		
		Aplicando esto a la función compuesta $(f \circ \gamma)(t) = f(\gamma(t))$:
		\[
		(f \circ \gamma)'(t) = \lim_{h \to 0} \frac{f(\gamma(t+h)) - f(\gamma(t))}{h}
		\]
		
		Sea $z = \gamma(t)$ y $z(h) = \gamma(t+h)$. Como $\gamma$ es diferenciable en $t$, es continua en dicho punto, por lo que cuando $h \to 0$, tenemos que $z(h) \to z$. 
		
		Además, como $f$ es diferenciable en el sentido complejo en $z$, por la definición de la derivada compleja, podemos escribir la función $f$ en un entorno de $z$ como:
		\[
		f(w) = f(z) + f'(z)(w - z) + R(w)(w - z)
		\]
		donde el término de error $R(w)$ satisface $\lim_{w \to z} R(w) = 0$.
		
		Sustituyendo $w = z(h) = \gamma(t+h)$:
		\[
		f(\gamma(t+h)) = f(\gamma(t)) + f'(\gamma(t)) \left( \gamma(t+h) - \gamma(t) \right) + R(\gamma(t+h)) \left( \gamma(t+h) - \gamma(t) \right)
		\]
		
		Sustituyendo esta expresión en el cociente incremental de la composición:
		\[
		\frac{f(\gamma(t+h)) - f(\gamma(t))}{h} = f'(\gamma(t)) \left( \frac{\gamma(t+h) - \gamma(t)}{h} \right) + R(\gamma(t+h)) \left( \frac{\gamma(t+h) - \gamma(t)}{h} \right)
		\]
		
		Tomamos el límite cuando $h \to 0$:
		\begin{enumerate}
			\item Para el primer término:
			\[
			\lim_{h \to 0} f'(\gamma(t)) \left( \frac{\gamma(t+h) - \gamma(t)}{h} \right) = f'(\gamma(t)) \cdot \gamma'(t)
			\]
			\item Para el segundo término: dado que $\lim_{h \to 0} \gamma(t+h) = \gamma(t) = z$, se sigue que $\lim_{h \to 0} R(\gamma(t+h)) = R(z) = 0$. Multiplicado por el límite finito $\gamma'(t)$, este término tiende a cero.
		\end{enumerate}
		
		Por lo tanto, se concluye que:
		\[
		(f \circ \gamma)'(t) = \{f'[\gamma(t)]\} \cdot \gamma'(t).
		\]
	\end{proof}
	
	\noindent \textbf{Ejercicio 48.} Para cada función $f : \Omega \to \mathbb{C}$ holomorfa en un conjunto abierto conexo $\Omega \subseteq \mathbb{C}$, pruebe las siguientes afirmaciones:
	
	\begin{enumerate}
		\item[(48.1)] Si $f'(z) = 0$ para todo $z \in \Omega$, entonces $f$ es constante.
		\item[(48.2)] Si existe $c \in \mathbb{C}$ tal que $f(z) = c \cdot \overline{f(z)}$ para todo $z \in \Omega$, entonces $f$ es constante.
		\item[(48.3)] Si $f(\Omega) \subseteq \mathbb{R}$, entonces $f$ es constante.
		\item[(48.4)] Si $\vert{}f\vert{}$ es constante, entonces $f$ es constante.
		\item[(48.5)] Si $g : \mathbb{C} \to \mathbb{C}$ es holomorfa y $g \circ f$ es constante, entonces $f$ o $g$ es constante.
		\item[(48.6)] Si $f_1, \ldots, f_N$ son holomorfas en $\Omega$, y si $\vert{}f_1\vert{}^2 + \cdots + \vert{}f_N\vert{}^2$ es constante, entonces cada $f_j$ es constante.
	\end{enumerate}
	
	\vspace{1em}
	\noindent \textbf{Demostraciones:}
	
	\begin{enumerate}
		\item[(48.1)] \textbf{Demostración:} Dado que $\Omega$ es un conjunto abierto y conexo, y $f'(z) = 0$ en todo $\Omega$, las ecuaciones de Cauchy-Riemann implican que las derivadas parciales de las partes real e imaginaria de $f$ son cero en todo $\Omega$. Por lo tanto, $f$ es localmente constante y, al ser $\Omega$ conexo, $f$ debe ser una función constante en todo $\Omega$. $\blacksquare$
		
		```
		\item[(48.2)] \textbf{Demostración:} Si $f(z) = c \cdot \overline{f(z)}$ para todo $z \in \Omega$, tomando el módulo en ambos lados obtenemos $\vert{}f(z)\vert{} = \vert{}c\vert{} \vert{}\overline{f(z)}\vert{} = \vert{}c\vert{} \vert{}f(z)\vert{}$. Si $f$ no es idénticamente cero, entonces $\vert{}c\vert{} = 1$. Expresando la condición analítica junto con las ecuaciones de Cauchy-Riemann, se deduce que las derivadas parciales de $f$ se anulan o que $\vert{}f\vert{}$ es constante, lo cual lleva a que $f$ sea una función constante. $\blacksquare$
		
		\item[(48.3)] \textbf{Demostración:} Si $f(\Omega) \subseteq \mathbb{R}$, la parte imaginaria de $f$ es idénticamente cero. Por las ecuaciones de Cauchy-Riemann, esto implica que $f'(z) = 0$ en todo $\Omega$. Aplicando el resultado del inciso (48.1), concluimos que $f$ es constante. $\blacksquare$
		
		\item[(48.4)] \textbf{Demostración:} Si $\vert{}f\vert{}$ es constante, sea $C$ su valor. Si $C = 0$, $f = 0$ trivialmente. Si $C > 0$, entonces $\vert{}f(z)\vert{}^2 = u^2 + v^2 = C^2$. Diferenciando esta expresión respecto a $x$ e $y$ y utilizando las ecuaciones de Cauchy-Riemann, se obtiene que las derivadas parciales de $u$ y $v$ son nulas, demostrando que $f$ es constante. $\blacksquare$
		
		\item[(48.5)] \textbf{Demostración:} Supongamos que $g(f(z)) = K$ para una constante $K$. Si $g$ no es constante, por el teorema de la aplicación abierta, la imagen de $g$ es un conjunto abierto, lo cual forzaría a que $f(\Omega)$ sea un conjunto unitario (es decir, $f$ es constante). Por lo tanto, necesariamente o bien $f$ es constante o bien $g$ es constante. $\blacksquare$
		
		\item[(48.6)] \begin{proof}
			Sean $f_1,\ldots,f_N$ funciones holomorfas en un abierto conexo
			$\Omega\subseteq\mathbb{C}$ y supongamos que
			\[
			|f_1(z)|^2+\cdots+|f_N(z)|^2=C
			\]
			para todo $z\in\Omega$, donde $C$ es una constante.
			
			Como $|f_j(z)|^2=f_j(z)\overline{f_j(z)}$, podemos escribir
			\[
			\sum_{j=1}^{N}f_j(z)\overline{f_j(z)}=C.
			\]
			
			Aplicando el laplaciano $\Delta=4\frac{\partial^2}
			{\partial z\,\partial\overline{z}}$ a ambos lados, obtenemos
			\[
			\Delta\left(\sum_{j=1}^{N}|f_j(z)|^2\right)=0.
			\]
			
			Por linealidad del laplaciano,
			\[
			\sum_{j=1}^{N}\Delta |f_j(z)|^2=0.
			\]
			
			Ahora bien, como cada $f_j$ es holomorfa, se cumple que
			\[
			\Delta |f_j(z)|^2=4|f_j'(z)|^2.
			\]
			
			Por consiguiente,
			\[
			4\sum_{j=1}^{N}|f_j'(z)|^2=0,
			\]
			y, dividiendo entre $4$, resulta
			\[
			\sum_{j=1}^{N}|f_j'(z)|^2=0.
			\]
			
			Como cada término $|f_j'(z)|^2$ es no negativo, necesariamente
			\[
			|f_j'(z)|^2=0
			\]
			para todo $j=1,\ldots,N$ y todo $z\in\Omega$. Por tanto,
			\[
			f_j'(z)=0
			\qquad \text{para todo }z\in\Omega.
			\]
			
			Finalmente, como $\Omega$ es conexo, por el resultado (48.1),
			cada función $f_j$ es constante en $\Omega$.
			
			En conclusión,
			\[
			\boxed{f_1,\ldots,f_N\text{ son constantes en }\Omega.}
			\]
		\end{proof}
		
		
	\end{enumerate}
	\noindent
	\textbf{Ejercicio 66.} Este ejercicio extiende el criterio de la razón de series reales a series generales.
	
	\vspace{1em}
	
	\begin{enumerate}[label=\textbf{(66.\arabic*)}, leftmargin=*]
		\item Para cada sucesión compleja $(w_n)$, demuestre que si 
		\[
		\limsup \frac{|w_{n+1}|}{|w_n|} < 1,
		\]
		entonces la serie compleja $\sum_{n=0}^{\infty} w_n$ converge.
		
		\item Inserte un ejemplo tal que 
		\[
		\limsup \frac{|w_{n+1}|}{|w_n|} = \infty
		\]
		y, aun así, la serie compleja $\sum_{n=0}^{\infty} w_n$ converja.
		
		\item Para cada sucesión compleja $(c_n)$, sea $r := 1 / (\limsup |c_{n+1}|/|c_n|)$. Muestre que la serie de potencias $\sum_{n=0}^{\infty} c_n z^n$ converge para cada $z \in D(0, r)$.
		
		\item[\textbf{$\dagger$ (66.4)}] Extienda los resultados anteriores a sucesiones y series en álgebras de Banach, por ejemplo, a matrices.
	\end{enumerate}
	
	\noindent \textbf{Ejercicio 67.} Pruebe que para cada serie de potencias $\sum_{n=0}^{\infty} c_n w^n$, la siguiente fórmula, atribuida a Jacques Hadamard, da el radio de convergencia $R$:
	\[
	R = \frac{1}{\limsup (\sqrt[n]{|c_n|})}.
	\]
	
	\begin{proof}
		Sea $R$ el radio de convergencia de la serie $\sum_{n=0}^{\infty} c_n w^n$. Por definición $R$ es el numero real en cual para $|z|<R$ la serie converge y para $|z|>R$ la serie diverge. De la definición, el numero real que cumple esta condición es unico, por tanto basta comprobar que $ \frac{1}{\limsup (\sqrt[n]{|c_n|})}$ cumple esta condición.
		
		Sea $|z|< \frac{1}{\limsup (\sqrt[n]{|c_n|})}$ entonces
		$$\limsup((\sqrt[n]{c_nz^n})= \limsup{\sqrt[n]{c_n}} z<1$$
		Por criterio de la raiz la serie converge.
		
		Analogamente, si suponemos que $|z|>1 \frac{1}{\limsup (\sqrt[n]{|c_n|})}$ tenemos que
		$$\limsup((\sqrt[n]{c_nz^n})= \limsup{\sqrt[n]{c_n}} z>1$$
		Por criterio de la raiz esta serie diverge. Así $\frac{1}{\limsup (\sqrt[n]{|c_n|})}$ es el radio de convergencia de la serie
		
	\end{proof}
	\vspace{1em}
	\noindent \textbf{Ejercicio 68.} Para cada sucesión de números complejos $(c_n)$, pruebe que si los límites $\lim(\sqrt[n]{|c_n|})$ y $\lim(|c_{n+1}| / |c_n|)$ existen (ninguno es infinito), entonces
	\[
	\lim(\sqrt[n]{|c_n|}) = \lim\left(\frac{|c_{n+1}|}{|c_n|}\right).
	\]
	\noindent \textbf{Demostración:}
	
	Sea $a_n = |c_n|$ una sucesión de números reales no negativos. Supongamos que el límite del cociente existe y es igual a $L \in \mathbb{R}$:
	\[
	\lim_{n \to \infty} \frac{a_{n+1}}{a_n} = L.
	\]
	
	Demostraremos que $\lim_{n \to \infty} \sqrt[n]{a_n} = L$. Para ello, analizaremos la definición de límite:
	
	Sea $\varepsilon > 0$ arbitrario. Como $\lim_{n \to \infty} \frac{a_{n+1}}{a_n} = L$, existe un entero $N \in \mathbb{N}$ tal que para todo $k \ge N$:
	\[
	L - \varepsilon < \frac{a_{k+1}}{a_k} < L + \varepsilon.
	\]
	
	Para cualquier $n > N$, podemos expresar $a_n$ como el producto telescópico:
	\[
	a_n = a_N \cdot \frac{a_{N+1}}{a_N} \cdot \frac{a_{N+2}}{a_{N+1}} \cdots \frac{a_n}{a_{n-1}}.
	\]
	
	Aplicando las acotaciones anteriores a cada uno de los $n - N$ factores del producto:
	\[
	a_N (L - \varepsilon)^{n-N} < a_n < a_N (L + \varepsilon)^{n-N}.
	\]
	
	Tomando la raíz $n$-ésima en toda la desigualdad:
	\[
	\sqrt[n]{a_N} (L - \varepsilon)^{\frac{n-N}{n}} < \sqrt[n]{a_n} < \sqrt[n]{a_N} (L + \varepsilon)^{\frac{n-N}{n}}.
	\]
	
	Reescribiendo los exponentes:
	\[
	\sqrt[n]{a_N} (L - \varepsilon)^{1 - \frac{N}{n}} < \sqrt[n]{a_n} < \sqrt[n]{a_N} (L + \varepsilon)^{1 - \frac{N}{n}}.
	\]
	
	Tomando el límite cuando $n \to \infty$ en ambos extremos, dado que $N$ y $a_N$ son fijos:
	\[
	\lim_{n \to \infty} \sqrt[n]{a_N} = 1 \quad \text{y} \quad \lim_{n \to \infty} \left(1 - \frac{N}{n}\right) = 1.
	\]
	
	Por lo tanto:
	\[
	L - \varepsilon \le \liminf_{n \to \infty} \sqrt[n]{a_n} \le \limsup_{n \to \infty} \sqrt[n]{a_n} \le L + \varepsilon.
	\]
	
	Como $\varepsilon > 0$ es totalmente arbitrario, esto implica que:
	\[
	\limsup_{n \to \infty} \sqrt[n]{a_n} = \liminf_{n \to \infty} \sqrt[n]{a_n} = L.
	\]
	
	En consecuencia, el límite existe y coincide con $L$:
	\[
	\lim_{n \to \infty} \sqrt[n]{|c_n|} = \lim_{n \to \infty} \frac{|c_{n+1}|}{|c_n|}. \quad \blacksquare
	\]
	
	\vspace{1em}
	\noindent \textbf{Ejercicio 69.} Para cada transformación fraccional lineal, determine su serie de Taylor (de potencias) centrada en cualquier punto donde esté definida.
	
	\noindent \textbf{Ejercicio 70.}
	
	\begin{enumerate}[label=\textbf{(70.\arabic*)}, leftmargin=*]
		\item Dé un ejemplo de una serie de potencias con un radio de convergencia igual a $1$ que \textit{converja} en todas partes sobre el círculo unitario $S^1 = \{z \in \mathbb{C} : |z| = 1\}$.
		
		\begin{proof}
			La serie $\sum \frac{z^n}{n^2}$. Su radio de convergencia es $1$ y converge en la frontera pues por criterio de comparación con $\sum 1/n^2$ que se sabe que es absolutamente convergente.
		\end{proof}
		
		\item Dé un ejemplo de una serie de potencias con un radio de convergencia igual a $1$ que \textit{diverja} en todas partes sobre el círculo unitario $S^1$.
		
		\item Dé un ejemplo de una serie de potencias con un radio de convergencia igual a $1$ que converja en algunos puntos pero diverja en otros puntos sobre el círculo unitario $S^1$.
		
		\begin{proof}
			La serie $\sum \frac{z^n}{n}$. Con $z=1$ diverge pero con $z=-1$ converge
		\end{proof}
		
		\item Determine la cardinalidad del conjunto $S$ de todas las series de potencias con un radio de convergencia igual a $1$.
		
		\item Deduzca que la siguiente afirmación es falsa: ``Dado cualquier subconjunto $E$ del círculo unitario $S^1$, existe un elemento de $S$ que converge en cada punto de $E$ y diverge en cada punto de su complemento $S^1 \setminus E$''.
	\end{enumerate}
	
	\section*{1.5 Exponential, trigonometric, and dilogarithm functions}
	
	\noindent \textbf{Exercise 71.} This exercise introduces the complex exponential function. With the notational convention $0^0 = 1$, define a function $\exp$ by:
	\begin{equation}\tag{9}
		\exp(z) := \sum_{n=0}^\infty \frac{z^n}{n!}
	\end{equation}
	
	\begin{enumerate}
		\item[(71.1)] Verify that the series in equation (9) converges in all of $\mathbb{C}$.
		\item[(71.2)] Verify that $\exp' = \exp$ also in the whole complex plane.
	\end{enumerate}
	
	\subsection*{Resolución}
	
	\subsubsection*{(71.1) Convergencia de la serie en todo $\mathbb{C}$}
	Para demostrar que la serie de potencias converge para todo $z \in \mathbb{C}$, utilizamos el \textbf{Criterio del Cociente (o de d'Alembert)} para series de términos complejos. 
	
	Sea $a_n = \frac{z^n}{n!}$. Evaluamos el límite del cociente de términos consecutivos en valor absoluto para $z \neq 0$:
	\[
	\lim_{n \to \infty} \left| \frac{a_{n+1}}{a_n} \right| = \lim_{n \to \infty} \left| \frac{\frac{z^{n+1}}{(n+1)!}}{\frac{z^n}{n!}} \right| = \lim_{n \to \infty} \left| \frac{z^{n+1} \cdot n!}{z^n \cdot (n+1)!} \right| = \lim_{n \to \infty} \frac{|z|}{n+1}
	\]
	
	Dado que para cualquier número complejo fijo $z$ se cumple que $\lim_{n \to \infty} \frac{|z|}{n+1} = 0$, este límite es estrictamente menor que $1$ para todo $z \in \mathbb{C}$. 
	
	Por el criterio del cociente, la serie converge absolutamente para todo $z \in \mathbb{C}$. Esto significa que el radio de convergencia es $R = \infty$.
	
	\subsubsection*{(71.2) Derivada de la función exponencial ($\exp' = \exp$)}
	Sabemos por la teoría de series de potencias en variable compleja que una serie de potencias puede derivarse término a término dentro de su disco de convergencia (en este caso, en todo $\mathbb{C}$). 
	
	Diferenciamos la serie respecto a $z$:
	\[
	\exp'(z) = \frac{d}{dz} \left( \sum_{n=0}^\infty \frac{z^n}{n!} \right) = \sum_{n=0}^\infty \frac{d}{dz} \left( \frac{z^n}{n!} \right)
	\]
	
	El término correspondiente a $n=0$ es una constante ($\frac{z^0}{0!} = 1$), por lo que su derivada es cero. Así, podemos iniciar la suma desde $n=1$:
	\[
	\exp'(z) = \sum_{n=1}^\infty \frac{n z^{n-1}}{n!}
	\]
	
	Simplificando el factorial $n! = n \cdot (n-1)!$:
	\[
	\exp'(z) = \sum_{n=1}^\infty \frac{n z^{n-1}}{n(n-1)!} = \sum_{n=1}^\infty \frac{z^{n-1}}{(n-1)!}
	\]
	
	Realizando un cambio de índice, sea $k = n - 1$. Cuando $n = 1$, $k = 0$; y cuando $n \to \infty$, $k \to \infty$:
	\[
	\exp'(z) = \sum_{k=0}^\infty \frac{z^k}{k!}
	\]
	
	Como esta última serie es exactamente la definición original de $\exp(z)$, concluimos que:
	\[
	\exp'(z) = \exp(z) \quad \text{para todo } z \in \mathbb{C}.
	\]
	
	\noindent \textbf{Ejercicio 72.} Este ejercicio establece algunas características de las funciones trigonométricas complejas. Con la convención de notación $0^0 = 1$, defínase:
	$$\cos(z) := \sum_{m=0}^\infty \frac{(-1)^m}{(2m)!} z^{2m},$$$$\sin(z) := \sum_{m=0}^\infty \frac{(-1)^m}{(2m+1)!} z^{2m+1}.$$\begin{enumerate}
		\item[(72.1)] Verificar que estas series para $\cos$ y $\sin$ convergen en todo $\mathbb{C}$.
		\item[(72.2)] Verificar que se cumplen las siguientes identidades para todo $z \in \mathbb{C}$:
		$$\begin{aligned} \exp(iz) &= \cos(z) + i \sin(z), \\ \cos(z) &= [\exp(iz) + \exp(-iz)]/2, \\ \sin(z) &= [\exp(iz) - \exp(-iz)]/(2i). \end{aligned}$$
		\item[(72.3)] Verificar que $\cos' = -\sin$ y $\sin' = \cos$.
	\end{enumerate}   \subsection*{Resolución}\subsubsection*{(72.1) Convergencia de las series en todo $\mathbb{C}$}\begin{proof}Para demostrar la convergencia en todo el plano complejo $\mathbb{C}$, aplicamos el Criterio del Cociente (o de d'Alembert) para series de potencias.\textbf{1. Para $\cos(z)$:} Sea $a_m = \frac{(-1)^m}{(2m)!} z^{2m}$. Evaluamos el límite del cociente de términos consecutivos en valor absoluto para $z \neq 0$:$$\lim_{m \to \infty} \left\vert{} \frac{a_{m+1}}{a_m} \right\vert{} = \lim_{m \to \infty} \left\vert{} \frac{\frac{(-1)^{m+1} z^{2(m+1)}}{(2(m+1))!}}{\frac{(-1)^m z^{2m}}{(2m)!}} \right\vert{} = \lim_{m \to \infty} \frac{\vert{}z\vert{}^2}{(2m+2)(2m+1)} = 0 < 1$$Como el límite es $0$ (estrictamente menor que $1$) para cualquier $z \in \mathbb{C}$, la serie converge absolutamente en todo el plano complejo (radio de convergencia $R = \infty$).\textbf{2. Para $\sin(z)$:} Sea $b_m = \frac{(-1)^m}{(2m+1)!} z^{2m+1}$. De manera análoga:$$\lim_{m \to \infty} \left\vert{} \frac{b_{m+1}}{b_m} \right\vert{} = \lim_{m \to \infty} \left\vert{} \frac{\frac{(-1)^{m+1} z^{2(m+1)+1}}{(2(m+1)+1)!}}{\frac{(-1)^m z^{2m+1}}{(2m+1)!}} \right\vert{} = \lim_{m \to \infty} \frac{\vert{}z\vert{}^2}{(2m+3)(2m+2)} = 0 < 1$$Por lo tanto, la serie de $\sin(z)$ también converge absolutamente para todo $z \in \mathbb{C}$.\end{proof}\subsubsection*{(72.2) Verificación de las identidades trigonométricas y exponenciales}\begin{proof}Partimos de la definición de la función exponencial compleja $\exp(iz) = \sum_{n=0}^\infty \frac{(iz)^n}{n!}$. Separamos la suma en índices pares ($n = 2m$) e impares ($n = 2m+1$):$$\exp(iz) = \sum_{m=0}^\infty \frac{(iz)^{2m}}{(2m)!} + \sum_{m=0}^\infty \frac{(iz)^{2m+1}}{(2m+1)!}$$Utilizando las potencias de la unidad imaginaria $i^{2m} = (i^2)^m = (-1)^m$ e $i^{2m+1} = i \cdot i^{2m} = i(-1)^m$:$$\exp(iz) = \sum_{m=0}^\infty \frac{(-1)^m z^{2m}}{(2m)!} + i \sum_{m=0}^\infty \frac{(-1)^m z^{2m+1}}{(2m+1)!} = \cos(z) + i \sin(z)$$Evaluando la expresión para $-z$:$$\exp(-iz) = \cos(-z) + i \sin(-z) = \cos(z) - i \sin(z)$$sumando y restando ambas ecuaciones obtenemos las expresiones algebraicas para $\cos(z)$ y $\sin(z)$:$$\cos(z) = \frac{\exp(iz) + \exp(-iz)}{2}, \quad \sin(z) = \frac{\exp(iz) - \exp(-iz)}{2i}$$\end{proof}\subsubsection*{(72.3) Derivadas de $\cos(z)$ y $\sin(z)$}\begin{proof}Derivamos término a término las series de potencias dentro de su dominio de convergencia en $\mathbb{C}$:\textbf{1. Derivada del coseno:}$$\cos'(z) = \frac{d}{dz} \left( \sum_{m=0}^\infty \frac{(-1)^m}{(2m)!} z^{2m} \right) = \sum_{m=1}^\infty \frac{(-1)^m 2m}{(2m)!} z^{2m-1}$$Simplificando $(2m)! = 2m(2m-1)!$:$$\cos'(z) = \sum_{m=1}^\infty \frac{(-1)^m}{(2m-1)!} z^{2m-1}$$Haciendo el cambio de índice $k = m - 1$ (de modo que $m = k + 1$ y $(-1)^{k+1} = -(-1)^k$):$$\cos'(z) = \sum_{k=0}^\infty \frac{-(-1)^k}{(2k+1)!} z^{2k+1} = - \sum_{k=0}^\infty \frac{(-1)^k}{(2k+1)!} z^{2k+1} = -\sin(z)$$\textbf{2. Derivada del seno:}$$\sin'(z) = \frac{d}{dz} \left( \sum_{m=0}^\infty \frac{(-1)^m}{(2m+1)!} z^{2m+1} \right) = \sum_{m=0}^\infty \frac{(-1)^m (2m+1)}{(2m+1)!} z^{2m}$$Simplificando $(2m+1)! = (2m+1)(2m)!$:$$\sin'(z) = \sum_{m=0}^\infty \frac{(-1)^m}{(2m)!} z^{2m} = \cos(z)$$\end{proof}	
	\noindent \textbf{Ejercicio 73.} Una de las definiciones del dilogaritmo $\operatorname{Li}_2$ es la serie
	\[
	\operatorname{Li}_2(z) := \sum_{n=1}^\infty \frac{z^n}{n^2}.
	\]
	
	\begin{enumerate}
		\item[(73.1)] Determinar su radio de convergencia, $R$.
		\item[(73.2)] Determinar si la serie converge en la clausura $\overline{D(0, R)}$.
		\item[(73.3)] Para cada real $\varepsilon > 0$, determinar un entero $n_\varepsilon$ tal que para todo entero $m \ge n_\varepsilon$ y para todo $z \in D(0, R)$
		\[
		\left| \sum_{n>m} \frac{z^n}{n^2} \right| < \varepsilon.
		\]
		\item[(73.4)] Mostrar que en el interior topológico del disco de convergencia, el dilogaritmo complejo satisface una ecuación diferencial ordinaria lineal de segundo orden con coeficientes racionales.
	\end{enumerate}
	
	\subsection*{Resolución}
	
	\subsubsection*{(73.1) Determinación del radio de convergencia $R$}
	\begin{proof}
		Utilizamos la fórmula de Cauchy-Hadamard para calcular el radio de convergencia $R$ de la serie de potencias $\sum_{n=1}^\infty a_n z^n$, donde los coeficientes son $a_n = \frac{1}{n^2}$:
		\[
		\frac{1}{R} = \limsup_{n \to \infty} \sqrt[n]{|a_n|} = \lim_{n \to \infty} \sqrt[n]{\frac{1}{n^2}} = \lim_{n \to \infty} \frac{1}{(\sqrt[n]{n})^2} = 1.
		\]
		Por lo tanto, el radio de convergencia es $R = 1$.
	\end{proof}
	
	\subsubsection*{(73.2) Convergencia en la clausura $\overline{D(0, R)}$}
	\begin{proof}
		La clausura del disco de convergencia es el disco cerrado $\overline{D(0, 1)} = \{z \in \mathbb{C} : |z| \le 1\}$. Para cualquier número complejo $z$ en esta clausura, tenemos la cota:
		\[
		\left| \frac{z^n}{n^2} \right| = \frac{|z|^n}{n^2} \le \frac{1}{n^2}.
		\]
		Dado que la serie numérica real $\sum_{n=1}^\infty \frac{1}{n^2}$ es convergente (por ser una p-serie con $p = 2 > 1$), el criterio de comparación de Weierstrass (M-test) garantiza que la serie de potencias converge absoluta y uniformemente en todo el disco cerrado $\overline{D(0, 1)}$. En conclusión, la serie sí converge en la clausura.
	\end{proof}
	
	\subsubsection*{(73.3) Estimación del resto de la serie}
	\begin{proof}
		Sea $\varepsilon > 0$ un número real arbitrario. Para cualquier $z$ en el disco de convergencia $D(0, R)$ (es decir, con $|z| \le 1$), acotamos el valor absoluto del resto mediante la desigualdad triangular:
		\[
		\left| \sum_{n>m} \frac{z^n}{n^2} \right| \le \sum_{n=m+1}^\infty \frac{|z|^n}{n^2} \le \sum_{n=m+1}^\infty \frac{1}{n^2}.
		\]
		Como la serie numérica $\sum_{n=1}^\infty \frac{1}{n^2}$ converge, la cola de la serie $\sum_{n=m+1}^\infty \frac{1}{n^2}$ tiende a cero cuando $m \to \infty$. Por consiguiente, existe un entero $n_\varepsilon$ tal que para todo $m \ge n_\varepsilon$ se cumple que:
		\[
		\sum_{n=m+1}^\infty \frac{1}{n^2} < \varepsilon,
		\]
		lo cual demuestra que $\left| \sum_{n>m} \frac{z^n}{n^2} \right| < \varepsilon$ para todo $z \in D(0, R)$.
	\end{proof}
	
	\subsubsection*{(73.4) Ecuación diferencial ordinaria para el dilogaritmo}
	\begin{proof}
		Definimos la función $f(z) = \operatorname{Li}_2(z) = \sum_{n=1}^\infty \frac{z^n}{n^2}$ en el interior del disco de convergencia $|z| < 1$. Derivando la serie término a término dentro del disco abierto, obtenemos la primera derivada:
		\[
		f'(z) = \sum_{n=1}^\infty \frac{n z^{n-1}}{n^2} = \sum_{n=1}^\infty \frac{z^{n-1}}{n}.
		\]
		Multiplicando ambos lados por $z$:
		\[
		z f'(z) = \sum_{n=1}^\infty \frac{z^n}{n} = -\ln(1-z).
		\]
		Derivamos esta expresión resultante respecto a la variable compleja $z$:
		\[
		\frac{d}{dz} [z f'(z)] = f'(z) + z f''(z) = \frac{1}{1-z}.
		\]
		Multiplicando toda la ecuación por $(1-z)$, obtenemos una ecuación diferencial ordinaria lineal de segundo orden con coeficientes polinómicos (racionales):
		\[
		(1-z) \left( z f''(z) + f'(z) \right) = 1,
		\]
		o de forma expandida:
		\[
		z(1-z) f''(z) + (1-z) f'(z) - 1 = 0,
		\]
		lo cual demuestra que el dilogaritmo complejo satisface esta EDO en el interior del disco de convergencia.
	\end{proof}
	

	Ejercicio 74. Pruebe que para cada función $f : \mathcal{D} \to \mathbb{C}$ continua en un subconjunto abierto $\mathcal{D} \subseteq \mathbb{C}$ y para cada curva diferenciable $\gamma : [0, 1] \to \mathcal{D}$,   $$\int_\gamma f(z) \, dz = \lim_{\max \vert{}t_{k+1} - t_k\vert{} \to 0} \sum_{k=0}^{N-1} f(z_k)(z_{k+1} - z_k)$$   
	para cada sucesión $(z_k)$ en la imagen de $\gamma$ tal que $z_k = \gamma(t_k)$ y $0 = t_0 < t_1 < \dots < t_{N-1} < t_N = 1$.   Resolución
	\begin{proof}
		Por definición, la integral de línea compleja de una función continua $f$ a lo largo de una curva diferenciable $\gamma : [0, 1] \to \mathcal{D}$ se define mediante la integral de Riemann en la variable real $t$:$$\int_\gamma f(z) \, dz = \int_0^1 f(\gamma(t)) \gamma'(t) \, dt$$Consideremos una partición del intervalo $[0, 1]$ dada por $0 = t_0 < t_1 < \dots < t_{N-1} < t_N = 1$, con una norma de partición que tiende a cero ($\max (t_{k+1} - t_k) \to 0$). Por la definición de la integral de Riemann como límite de sumas de Riemann, podemos expresar la integral anterior como:$$\int_0^1 f(\gamma(t)) \gamma'(t) \, dt = \lim_{\max \Delta t_k \to 0} \sum_{k=0}^{N-1} f(\gamma(t_k)) \gamma'(t_k) \Delta t_k$$donde $\Delta t_k = t_{k+1} - t_k$.Dado que $\gamma$ es una curva diferenciable, el incremento de la curva en cada subintervalo se relaciona con su derivada mediante la aproximación lineal:$$\gamma(t_{k+1}) - \gamma(t_k) \approx \gamma'(t_k)(t_{k+1} - t_k) = \gamma'(t_k) \Delta t_k$$Si denotamos $z_k = \gamma(t_k)$, entonces los puntos de la partición sobre la imagen de $\gamma$ cumplen que $z_{k+1} - z_k \approx \gamma'(t_k) \Delta t_k$. Sustituyendo esta relación en la suma de Riemann, el límite de la sumatoria se transforma en:$$\int_\gamma f(z) \, dz = \lim_{\max \vert{}t_{k+1} - t_k\vert{} \to 0} \sum_{k=0}^{N-1} f(z_k)(z_{k+1} - z_k)$$lo cual completa la demostración.
	\end{proof}
	
	\noindent \textbf{Ejercicio 75.} Pruebe que si $f : \mathcal{D} \to \mathbb{C}$ tiene una derivada compleja que es continua en el conjunto abierto $\mathcal{D} \subseteq \mathbb{C}$, y si $\gamma : [0, 1] \to \mathcal{D}$ es una curva diferenciable desde $z_0 := \gamma(0)$ hasta $z_1 := \gamma(1)$, entonces
	\[
	f(z_1) - f(z_0) = \int_{\gamma} f'(\zeta) \, d\zeta.
	\]
	
	\begin{proof}
		Aplicar directamente el lema del capitulo 1
	\end{proof}
	
	\noindent \textbf{Ejercicio 76.} Pruebe que si $f : \mathcal{D} \to \mathbb{C}$ tiene una derivada compleja $f'$ que es continua en el conjunto abierto $\mathcal{D} \subseteq \mathbb{C}$, y si $\gamma : [0, 1] \to \mathbb{C}$ es una curva diferenciable desde $z_0 := \gamma(0)$ hasta $z_1 := \gamma(1)$, entonces
	\[
	|f(z_1) - f(z_0)| \leq \text{Length}(\gamma) \cdot \max\{|f'(z)| : z \in \gamma([0, 1])\}.
	\]
	
	\vspace{1em}
	\noindent \textbf{Demostración}
	Por el teorema fundamental del cálculo para integrales de línea (ejercicio anterior), sabemos que:
	\[
	f(z_1) - f(z_0) = \int_{\gamma} f'(\zeta) \, d\zeta.
	\]
	Usando la definición de la integral de línea y la desigualdad estándar para integrales de funciones continuas, tenemos que:
	\[
	\left| \int_{\gamma} f'(\zeta) \, d\zeta \right| \leq \int_0^1 |f'(\gamma(t))| |\gamma'(t)| \, dt.
	\]
	Sea $M = \max\{|f'(z)| : z \in \gamma([0, 1])\}$. Podemos acotar el integrando:
	\[
	\int_0^1 |f'(\gamma(t))| |\gamma'(t)| \, dt \leq M \int_0^1 |\gamma'(t)| \, dt.
	\]
	Dado que la longitud de la curva $\gamma$ está definida como $\text{Length}(\gamma) = \int_0^1 |\gamma'(t)| \, dt$, se concluye directamente que:
	\[
	|f(z_1) - f(z_0)| \leq \text{Length}(\gamma) \cdot M. \quad \blacksquare
	\]
	
	\noindent \textbf{Ejercicio 77.} Pruebe que si una curva diferenciable $\gamma : [0, 1] \to \mathbb{C}$ parametriza en sentido antihorario la frontera $\partial\Omega$ de un conjunto abierto $\Omega \subset \mathbb{C}$, y si $A(\Omega)$ denota el área de $\Omega$, entonces bajo condiciones adecuadas sobre $\Omega$[cite: 6]:
	\[
	A(\Omega) = \frac{1}{2i} \oint_{\partial\Omega} \bar{z} \, dz. 
	\]
	
	\noindent \textbf{Demostración del Ejercicio 77:}
	
	Para demostrar la fórmula del área mediante la integral de línea en el plano complejo, podemos utilizar el Teorema de Green. 
	
	Sea $z = x + iy$, de modo que la integral de línea se puede expresar en términos de las coordenadas cartesianas $x$ e $y$. El diferencial complejo es:
	\[
	dz = dx + i\,dy
	\]
	y el conjugado complejo es:
	\[
	\bar{z} = x - iy.
	\]
	
	Sustituyendo estas expresiones en la integral de línea sobre la frontera $\partial\Omega$:
	\[
	\oint_{\partial\Omega} \bar{z} \, dz = \oint_{\partial\Omega} (x - iy)(dx + i\,dy)
	\]
	Expandiendo el producto:
	\[
	= \oint_{\partial\Omega} \left[ x\,dx + y\,dy + i(x\,dy - y\,dx) \right]
	\]
	Podemos separar esta integral en su parte real y su parte imaginaria:
	\[
	= \oint_{\partial\Omega} (x\,dx + y\,dy) + i \oint_{\partial\Omega} (x\,dy - y\,dx).
	\]
	
	Por el Teorema de Green en el plano, para un campo vectorial $\mathbf{F} = (P, Q)$ integrado en sentido antihorario, tenemos:
	\[
	\oint_{\partial\Omega} (P\,dx + Q\,dy) = \iint_n \left( \frac{\partial Q}{\partial x} - \frac{\partial P}{\partial y} \right) dA.
	\]
	
	1. Para la primera integral, $P = x$ y $Q = y$:
	\[
	\frac{\partial Q}{\partial x} - \frac{\partial P}{\partial y} = \frac{\partial}{\partial x}(y) - \frac{\partial}{\partial y}(x) = 0 - 0 = 0.
	\]
	Por lo tanto, $\oint_{\partial\Omega} (x\,dx + y\,dy) = 0$.
	
	2. Para la segunda integral, $P = -y$ y $Q = x$:
	\[
	\frac{\partial Q}{\partial x} - \frac{\partial P}{\partial y} = \frac{\partial}{\partial x}(x) - \frac{\partial}{\partial y}(-y) = 1 - (-1) = 2.
	\]
	Aplicando el Teorema de Green:
	\[
	\oint_{\partial\Omega} (x\,dy - y\,dx) = \iint_{\Omega} 2 \, dA = 2 A(\Omega).
	\]
	
	Sustituyendo estos resultados de vuelta en nuestra expresión inicial para la integral compleja:
	\[
	\oint_{\partial\Omega} \bar{z} \, dz = 0 + i \left( 2 A(\Omega) \right) = 2i \, A(\Omega).
	\]
	
	Finalmente, despejando el área $A(\Omega)$, obtenemos:
	\[
	A(\Omega) = \frac{1}{2i} \oint_{\partial\Omega} \bar{z} \, dz. \quad \blacksquare
	\]
	
	\noindent \textbf{Ejercicio 78.} Para cada rectángulo cerrado $R$ en un subconjunto abierto $\Omega \subseteq \mathbb{R}^2$, y para cada función continuamente diferenciable $f : \Omega \to \mathbb{R}^2$, aplique la siguiente versión del Teorema de Stokes del texto:
	\[
	\iint_{R} \frac{\partial f}{\partial \bar{z}} \, dx \, dy = \frac{1}{2i} \oint_{\partial R} f(z) \, dz
	\]
	para diseñar una demostración de que si $f \in C^1(\Omega)$ y $\oint_{\partial R} f(z) \, dz = 0$ para todo rectángulo cerrado $R$ en $\Omega$, entonces $f$ tiene una derivada compleja en todas partes en $\Omega$.
	
	\vspace{1em}
	\noindent \textbf{Demostración:}
	Por hipótesis, tenemos que para cada rectángulo cerrado $R \subseteq \Omega$, se cumple que:
	\[
	\oint_{\partial R} f(z) \, dz = 0.
	\]
	
	Utilizando la versión del Teorema de Stokes proporcionada en el enunciado, relacionamos la integral de línea sobre la frontera del rectángulo con una integral doble sobre $R$[cite: 7]:
	\[
	\frac{1}{2i} \oint_{\partial R} f(z) \, dz = \iint_{R} \frac{\partial f}{\partial \bar{z}} \, dx \, dy.
	\]
	
	Dado que la integral de línea es cero para cualquier rectángulo cerrado $R$, la integral doble debe ser cero para todo rectángulo $R \subset \Omega$:
	\[
	\iint_{R} \frac{\partial f}{\partial \bar{z}} \, dx \, dy = 0.
	\]
	
	Como la función derivada parcial $\frac{\partial f}{\partial \bar{z}}$ es continua en $\Omega$ (dado que $f \in C^1(\Omega)$), si la integral de esta función sobre cualquier rectángulo arbitrario es cero, el integrando mismo debe ser idénticamente cero en todo punto de $\Omega$. Es decir:
	\[
	\frac{\partial f}{\partial \bar{z}} = 0 \quad \text{en todo punto de } \Omega.
	\]
	
	Recordemos que la condición $\frac{\partial f}{\partial \bar{z}} = 0$ es equivalente a las ecuaciones de Cauchy-Riemann en forma compleja. Dado que $f \in C^1(\Omega)$ y sus derivadas parciales satisfacen las ecuaciones de Cauchy-Riemann en todo $\Omega$, se concluye que $f$ es una función holomorfa (es decir, tiene una derivada compleja en todas partes) en $\Omega$. $\blacksquare$
	
	\noindent \textbf{Definición 22.} Para cada número complejo $c \in \mathbb{C}$ y para cada curva cerrada diferenciable $\gamma : [a, b] \to \mathbb{C} \setminus \{c\}$ que no pasa por el punto $c$, defina el \textbf{número de índice} o \textbf{número de vueltas} $n(\gamma, c)$ de la curva $\gamma$ con respecto al punto $c$ como
	\[
	n(\gamma, c) := \frac{1}{2\pi i} \oint_{\gamma} \frac{1}{z - c} \, dz. \quad \square
	\]
	
	\vspace{1.5em}
	\noindent \textbf{Definición 23.} Sea $\sigma : [a, b] \to S^1 \subset \mathbb{C}$ una parametrización diferenciable cualquiera del círculo unitario en sentido antihorario. Para cada función compleja diferenciable $f : \Omega \to \mathbb{C}$ definida en una vecindad $\Omega$ del círculo unitario y tal que $0 \notin f(S^1)$, defina el \textbf{grado} $\text{Deg}(f, 0)$ de la función $f$ con respecto al origen como
	\[
	\text{Deg}(f, 0) := n(f \circ \sigma, 0).
	\]
	Por lo tanto, $\text{Deg}(f, 0)$ es el número de vueltas de la imagen de $S^1$ bajo $f$. $\square$
	
	\noindent \textbf{Ejercicio 80.} Pruebe que para cada entero $n \in \mathbb{Z}$, la función $z \mapsto z^{-n}$ tiene una primitiva en $\mathbb{C}^*$ si y solo si $n \neq 1$.
	
	\vspace{1em}
	\noindent \textbf{Demostración:}
	
	Una función continua $f : \Omega \to \mathbb{C}$ tiene una primitiva en un conjunto abierto y conexo $\Omega$ si y solo si su integral a lo largo de cualquier curva cerrada rectificable en $\Omega$ es cero, o equivalentemente, si la integral a lo largo de cualquier camino depende únicamente de sus puntos inicial y final. 
	
	Consideremos la función $f(z) = z^{-n}$ en el dominio abierto $\mathbb{C}^* = \mathbb{C} \setminus \{0\}$. Analizamos ambos casos:
	
	\begin{enumerate}
		\item[\textbf{1.}] \textbf{Si $n \neq 1$:}
		Podemos construir explícitamente una primitiva para $f(z)$. Definimos la función candidata:
		\[
		F(z) = \frac{z^{1-n}}{1-n}.
		\]
		Dado que $n \in \mathbb{Z}$ y $n \neq 1$, el exponente $1-n$ es un entero distinto de cero, por lo que $F(z)$ está bien definida y es holomorfa en todo $\mathbb{C}^*$. Al derivar $F(z)$ con respecto a $z$, obtenemos por la regla de la potencia:
		\[
		F'(z) = \frac{1}{1-n} \cdot (1-n) z^{1-n-1} = z^{-n} = f(z).
		\]
		Por lo tanto, $f(z)$ admite una primitiva en $\mathbb{C}^*$ cuando $n \neq 1$.
		
		\item[\textbf{2.}] \textbf{Si $n = 1$:}
		La función toma la forma $f(z) = z^{-1} = \frac{1}{z}$. 
		Supongamos por contradicción que existe una primitiva holomorfa $F(z)$ tal que $F'(z) = \frac{1}{z}$ en todo $\mathbb{C}^*$. 
		Consideremos la circunferencia unitaria recorrida en sentido antihorario, parametrizada por $\gamma(t) = e^{it}$ para $t \in [0, 2\pi]$, la cual se encuentra enteramente en $\mathbb{C}^*$. 
		Por el teorema fundamental del cálculo para integrales de línea, la integral de una función que posee una primitiva a lo largo de cualquier curva cerrada debe ser cero:
		\[
		\oint_{\gamma} \frac{1}{z} \, dz = F(\gamma(2\pi)) - F(\gamma(0)) = F(1) - F(1) = 0.
		\]
		Sin embargo, al calcular directamente esta integral, obtenemos:
		\[
		\oint_{\gamma} \frac{1}{z} \, dz = \int_0^{2\pi} \frac{1}{e^{it}} \cdot i e^{it} \, dt = i \int_0^{2\pi} 1 \, dt = 2\pi i \neq 0.
		\]
		Esto genera una contradicción directa ($2\pi i = 0$), lo cual demuestra que no puede existir ninguna primitiva para $z^{-1}$ en todo $\mathbb{C}^*$.
	\end{enumerate}
	
	De este modo, se concluye que la función $z \mapsto z^{-n}$ tiene una primitiva en $\mathbb{C}^*$ si y solo si $n \neq 1$. $\blacksquare$
	
	\noindent \textbf{Ejercicio 83.} Pruebe que la representación integral de Cauchy a lo largo del círculo unitario $S^1(0, 1)$,
	\[
	f(z) = \frac{1}{2\pi i} \oint_{S^1(0, 1)} \frac{f(w)}{w - z} \, dw,
	\]
	también se cumple para cada $z \in D(0, 1)$ y cada función $f : \overline{D(0, 1)} \to \mathbb{C}$ que es diferenciable en el sentido complejo en el disco unitario abierto $D(0, 1)$, y continua —pero no necesariamente diferenciable— en el disco unitario cerrado $\overline{D(0, 1)}$.
	
	\vspace{1em}
	\noindent \textbf{Demostración:}
	
	Para demostrar la fórmula de Cauchy bajo la suposición de que $f$ solo es continua en el disco cerrado $\overline{D(0, 1)}$ y holomorfa en el disco abierto $D(0, 1)$, podemos utilizar un argumento de aproximación mediante homotopía o contracción de círculos concéntricos.
	
	Sea $z \in D(0, 1)$ fijo. Dado que $D(0, 1)$ es un conjunto abierto, existe un radio $r$ tal que $|z| < r < 1$. 
	Consideremos el círculo de radio $r$ centrado en el origen, denotado por $S^1(0, r)$, orientado en sentido antihorario. 
	
	Dado que $f$ es holomorfa en $D(0, 1)$, podemos aplicar la fórmula integral de Cauchy estándar sobre el círculo interior $S^1(0, r)$ que encierra al punto $z$:
	\[
	f(z) = \frac{1}{2\pi i} \oint_{S^1(0, r)} \frac{f(w)}{w - z} \, dw.
	\]
	
	Queremos demostrar que esta integral sobre $S^1(0, r)$ es igual a la misma integral evaluada sobre el círculo unitario $S^1(0, 1)$. 
	Para ello, consideramos la corona circular (anillo) definida por $r \leq |w| \leq 1$. En esta región cerrada y acotada, la función $w \mapsto \frac{f(w)}{w - z}$ es continua (ya que $f$ es continua en $\overline{D(0, 1)}$ y el denominador no se anula porque $|w| \ge r > |z|$).
	
	Aplicando el teorema de deformación de contornos (o el teorema integral de Cauchy para dominios multipunturados), las integrales de una función holomorfa sobre curvas homótopas en un dominio son iguales. En particular, a medida que hacemos tender $r$ hacia $1$ ($r \to 1^-$), el círculo $S^1(0, r)$ se aproxima continuamente al círculo unitario $S^1(0, 1)$.
	
	Debido a la continuidad uniforme de $f$ en la cerradura del disco $\overline{D(0, 1)}$, el valor de la integral sobre los círculos concéntricos varía de manera continua respecto al radio $r$. Tomando el límite cuando $rro \to 1^-$, la integral sobre el círculo $S^1(0, r)$ converge a la integral sobre el círculo unitario $S^1(0, 1)$:
	\[
	\lim_{r \to 1^-} \frac{1}{2\pi i} \oint_{S^1(0, r)} \frac{f(w)}{w - z} \, dw = \frac{1}{2\pi i} \oint_{S^1(0, 1)} \frac{f(w)}{w - z} \, dw.
	\]
	
	Dado que el valor de la integral en el lado izquierdo es constante e igual a $f(z)$ para cualquier $r$ tal que $|z| < r < 1$, se concluye que:
	\[
	f(z) = \frac{1}{2\pi i} \oint_{S^1(0, 1)} \frac{f(w)}{w - z} \, dw. \quad \blacksquare
	\]
	
	\noindent \textbf{(Recíproco del ejercicio 71.)} Pruebe que si una función, denotada por $\exp$, es holomorfa y satisface la ecuación diferencial $\exp' = \exp$ en una vecindad de $0$, con valor inicial $\exp(0) = 1$, entonces, con la convención notacional $0^0 = 1$,
	\[
	\exp(z) = \sum_{n=0}^{\infty} \frac{z^n}{n!}.
	\]
	
	\vspace{1em}
	\noindent \textbf{Demostración:}
	
	Como la función $\exp$ es holomorfa en una vecindad del origen $D(0, R)$, admite un desarrollo en serie de Taylor alrededor de $z = 0$:
	\[
	\exp(z) = \sum_{n=0}^{\infty} a_n z^n, \quad \text{donde } a_n = \frac{\exp^{(n)}(0)}{n!}.
	\]
	
	A partir de la ecuación diferencial dada, $\exp'(z) = \exp(z)$, podemos calcular inductivamente las derivadas de orden superior de la función $\exp$:
	\begin{align*}
		\exp^{(1)}(z) &= \exp'(z) = \exp(z), \\
		\exp^{(2)}(z) &= (\exp'(z))' = \exp'(z) = \exp(z), \\
		&\;\vdots \\
		\exp^{(n)}(z) &= \exp(z) \quad \text{para todo } n \geq 0.
	\end{align*}
	
	Evaluando cada derivada en el punto $z = 0$ y utilizando la condición inicial $\exp(0) = 1$:
	\[
	\exp^{(n)}(0) = \exp(0) = 1 \quad \text{para todo } n \geq 0.
	\]
	
	Sustituyendo el valor de las derivadas en la fórmula de los coeficientes de Taylor $a_n$:
	\[
	a_n = \frac{1}{n!} \quad \text{para todo } n \geq 0.
	\]
	
	Por lo tanto, la representación en serie de Taylor de la función $\exp(z)$ alrededor del origen resulta ser:
	\[
	\exp(z) = \sum_{n=0}^{\infty} \frac{1}{n!} z^n = \sum_{n=0}^{\infty} \frac{z^n}{n!}. \quad \blacksquare
	\]
	
	\noindent \textbf{Ejercicio 103.} Sea $\gamma : [0, 1] \to \mathbb{C}$ un camino diferenciable a trozos tal que $\gamma(0) = 2$ y $\gamma^{-1}([0, \infty)) = \{0\}$[cite: 14]. Pruebe que
	\[
	0 < \mathfrak{Im} \left[ \int_0^1 \frac{\gamma'}{\gamma}(t) \, dt \right] < 2\pi. \text{[cite: 14]}
	\]
	
	\vspace{1em}
	\noindent \textbf{Demostración:}
	
	Consideremos la integral sobre el camino $\gamma$:
	\[
	I = \int_{\gamma} \frac{dz}{z} = \int_0^1 \frac{\gamma'(t)}{\gamma(t)} \, dt.
	\]
	
	Dado que la condición $\gamma^{-1}([0, \infty)) = \{0\}$ establece que el camino solo toca el semieje real no negativo en el punto inicial $t = 0$ con $\gamma(0) = 2$[cite: 14], el camino para $t \in (0, 1]$ permanece completamente contenido en el plano complejo con un corte de rama a lo largo de $[0, \infty)$, es decir, en $\mathbb{C} \setminus [0, \infty)$.
	
	En el dominio simplemente conexo $D = \mathbb{C} \setminus [0, \infty)$, podemos definir una rama continua y holomorfa del logaritmo complejo, denotada por $\log_{\theta}(z)$, donde la parte imaginaria (el argumento principal para esta corte de rama) está bien restringida al intervalo continuo de argumentos $(0, 2\pi)$.
	
	Para cualquier $t \in (0, 1]$, tenemos una primitiva para $\frac{1}{z}$ en este dominio acotado dada por $\log(z)$. Aplicando el Teorema Fundamental del Cálculo para integrales de línea en el intervalo $[\varepsilon, 1]$ con $\varepsilon > 0$:
	\[
	\int_{\varepsilon}^1 \frac{\gamma'(t)}{\gamma(t)} \, dt = \log(\gamma(1)) - \log(\gamma(\varepsilon)).
	\]
	
	Tomando el límite cuando $\varepsilon \to 0^+$, la integral se evalúa como la variación continua del logaritmo a lo largo de $\gamma$:
	\[
	\int_0^1 \frac{\gamma'(t)}{\gamma(t)} \, dt = \log(\gamma(1)) - \log(\gamma(0)).
	\]
	
	Al tomar la parte imaginaria en ambos lados:
	\[
	\mathfrak{Im} \left[ \int_0^1 \frac{\gamma'(t)}{\gamma(t)} \, dt \right] = \arg(\gamma(1)) - \arg(\gamma(0)).
	\]
	
	Como $\gamma(0) = 2 \in (0, \infty)$[cite: 14], el argumento inicial cuando nos aproximamos desde el corte de rama positivo se toma continuo dentro del intervalo de la rama $(0, 2\pi)$. Además, dado que $\gamma(t) \notin [0, \infty)$ para todo $t \in (0, 1]$[cite: 14], la función continua del argumento $t \mapsto \arg(\gamma(t))$ nunca cruza el eje real positivo en $t \in (0, 1]$.
	
	Por la definición continua de la parte imaginaria del logaritmo en la rama elegida sobre $\mathbb{C} \setminus [0, \infty)$, la variación del argumento $\arg(\gamma(1))$ satisface estrictamente:
	\[
	0 < \arg(\gamma(1)) < 2\pi.
	\]
	
	Por consiguiente, la parte imaginaria de la integral satisface la desigualdad deseada[cite: 14]:
	\[
	0 < \mathfrak{Im} \left[ \int_0^1 \frac{\gamma'}{\gamma}(t) \, dt \right] < 2\pi. \quad \blacksquare
	\]
\end{document}